% HAND-WRITTEN fragment. Design requirements: none of the five elements below
% exists in the current schematics. Evidence, so this can be re-checked:
%   skill/castalia_anatop_build.il buildChannel() wires the transimpedance
%   output net afe_out straight to the converter's Analog_Input, and brings
%   CE/RE/WE out of the channel unswitched; SIM_STATUS.md 3.1 confirms the same
%   for the pixel's schematic_adc view (converter input directly connected).
% The control bits land in the existing per-site switch-matrix word, so no
% digital-contract change is implied (include/CqAnalog.tex, offset 0x0C).
% Library cell names are kept out of the rendered prose.

\subsubsection{Measurement paths the procedure requires}
\label{sss:calibration-paths}

Every step below is a measurement the channel makes on itself through its own
converter. Five elements carry those measurements and \textbf{none of them is in
the current schematics}: the channel as drawn wires the transimpedance output
straight to the converter input and brings all three electrode nodes out
unswitched. They are design requirements, listed with what each one is for and
what it costs.

\begin{enumerate}[label=C\arabic*., leftmargin=*, itemsep=2pt]
\item \textbf{Series isolation switches in the counter-, reference- and
  working-electrode paths.} Without them no zero can be measured against
  anything but the electrode itself. The reference switch carries no current and
  the counter switch sits inside A1's loop, so only the working-electrode switch
  has a resistance specification: it carries the full cell current into the
  virtual ground, so its drop displaces the working electrode by
  \(i\!\cdot\!R_{\mathrm{on}}\) --- \SI{3.3}{\milli\volt} at
  \SI{33}{\micro\ampere} through \SI{100}{\ohm}. The current is the measured
  quantity, so the drop is computable and the requirement is that
  \(R_{\mathrm{on}}\) be \emph{stable and stored}, not that it be small; a
  quarter-LSB uncorrected drop at full scale would need \SI{12}{\ohm}.
\item \textbf{A unity-gain loopback switch around A1}, tying the
  counter-electrode node back to the amplifier's feedback input with the
  reference switch open. This is the only way to see A1's offset: it puts the
  amplifier in unity gain at the same input common mode it operates at, and
  leaves a low-impedance node to sample. Buffering the reference electrode into
  the converter instead is rejected --- the converter's sampling capacitor would
  draw charge from the reference electrode, which is electrochemically
  forbidden.
\item \textbf{A converter input multiplexer, at least 4:1}: the transimpedance
  output (normal operation), the common-mode rail, this channel's reference DAC
  output, and the counter-electrode node. The first two are the zero
  measurements, the third carries the reference DAC's own linearity, the fourth
  carries A1.
\item \textbf{An on-chip gain reference resistor}, switched between the
  counter-electrode node and the working-electrode summing node, and built from
  the same unit device as the feedback ladder. With the electrodes isolated and
  the loopback closed, the voltage across it is \(v_{\mathrm{ref}} -
  V_{\mathrm{CM}}\) and the converter reads \(V_{\mathrm{CM}} +
  (v_{\mathrm{ref}} - V_{\mathrm{CM}})\,R_f/R_{\mathrm{cal}}\). The sheet
  resistance cancels in that ratio, so the channel measures its own
  transimpedance to matching accuracy --- \SIrange{0.04}{0.36}{\percent}
  (Table~\ref{tab:tiarprog-mc}) --- with no external current and no reference of
  any kind. Take it as a two-point slope, so that both amplifier offsets cancel.
  Size the resistor near the middle of the ladder and choose the drive amplitude
  per gain code, or the high-gain codes rail.
\item \textbf{Control bits, eight to ten per channel.} They fit the per-site
  switch-matrix and analog-multiplexer word already reserved in the register map
  (Section~\ref{ss:cqanalog-regs}), so the digital contract does not change.
\end{enumerate}

\paragraph{Interface to the analog test multiplexer.} Two of the constants are
absolute and cannot be measured on-die, so two internal nodes must reach a pad
through the analog debug multiplexer. Both are needed only at production test,
at DC, and both tolerate a buffer in the path provided its offset is inside a
converter LSB.

\begin{itemize}[leftmargin=*, itemsep=2pt]
\item \textbf{Reference DAC output, per channel} --- the potential-axis absolute
  anchor. A voltmeter at three codes fixes the DAC's offset and step size in
  volts, which no on-chip measurement can supply.
\item \textbf{Common-mode rail} --- shared, and the zero of both axes. Measured
  once alongside the reference DAC so the two are on the same external scale.
\item \textbf{Counter-electrode node, per channel} --- optional. The loopback of
  C2 makes A1's offset readable on-die; a pad on the same node makes it directly
  observable during bring-up and is the cheapest way to prove the loopback
  itself.
\item \textbf{Transimpedance output, per channel} --- optional, and the same
  node the compensation study wants brought out. It turns the converter's own
  gain and linearity into a measurable quantity instead of an assumed one, which
  is the one place this chapter has no data at all.
\end{itemize}

\noindent The current-axis anchor needs no multiplexer: the tester forces its
reference current into the working-electrode pad, which is already bonded.
